\chapter{ECG}
\index{ECG}

\section*{Definition and Overview}
An electrocardiogram, known as an ECG or EKG, is a simple, painless test that records the electrical activity of the heart. Every heartbeat is triggered by a small electrical impulse that spreads through the heart muscle, and the ECG traces this activity as a series of waves, allowing doctors to assess heart rate and rhythm, detect damage from heart attacks, and identify other heart conditions. The ECG is one of the most common and useful tests in medicine, performed in doctors' surgeries, hospitals, ambulances, and emergency departments.

\section*{Epidemiology}
ECGs are performed millions of times each year in many countries. They are used to investigate chest pain, palpitations, breathlessness, and fainting, to screen people with risk factors or before surgery, and to monitor people with known heart disease. Because heart disease is the leading cause of death in many countries, the ECG plays a central role in the rapid diagnosis and treatment of heart attacks and arrhythmias.

\section*{Aetiology and Pathophysiology}
The heart's electrical system drives every beat.
\begin{itemize}
    \item \textbf{The pacemaker:} the sinoatrial node in the right atrium generates the electrical impulse that starts each heartbeat
    \item \textbf{Conduction:} the impulse spreads across the atria, causing them to contract, then passes through the atrioventricular node and specialised pathways to the ventricles
    \item \textbf{The ECG trace:} electrodes on the skin detect these electrical signals, producing waves labelled P, QRS, and T, which correspond to atrial contraction, ventricular contraction, and ventricular recovery
    \item \textbf{What ECGs detect:} heart rate and rhythm, conduction problems, heart attack damage, enlargement of the heart chambers, and the effects of some medicines and electrolyte disturbances
    \item \textbf{Limitations:} a normal ECG does not exclude all heart problems, and abnormalities may appear only at rest, during exercise, or at certain times
\end{itemize}

\section*{Clinical Features}
An ECG may be requested for a range of symptoms and situations.
\begin{itemize}
    \item Chest pain, pressure, or discomfort
    \item Palpitations or an irregular heartbeat
    \item Breathlessness, dizziness, or fainting
    \item A routine check for people with high blood pressure, diabetes, or heart disease risk factors
    \item Pre-operative assessment
    \item Monitoring after a heart attack or heart surgery
    \item Checking the effect of heart medicines
    \item Exercise stress testing, where the ECG is recorded during physical activity
\end{itemize}

\section*{Differential Diagnosis}
The ECG is a diagnostic tool, not a diagnosis in itself, and its interpretation depends on the clinical context.
\begin{itemize}
    \item A resting ECG may be normal despite significant coronary artery disease
    \item Some ECG changes are harmless variants, while others signal serious disease
    \item Comparison with previous ECGs is often essential for interpretation
    \item The ECG result is always combined with symptoms, examination, and blood tests
\end{itemize}

\section*{When to Seek Medical Care}
An ECG is performed when a doctor considers it necessary. Chest pain, palpitations, breathlessness, or fainting should always be discussed with a doctor, who can decide whether an ECG is needed.

\textbf{Contact your local emergency services for:}
\begin{itemize}
    \item Chest pain that is severe, crushing, or lasting more than a few minutes, especially with sweating, breathlessness, or nausea
    \item Fainting or collapse
    \item Palpitations with dizziness, chest pain, or breathlessness
    \item Signs of a stroke, including facial drooping, weakness, or difficulty speaking
\end{itemize}

\section*{Diagnosis and Investigations}
Performing an ECG is quick and straightforward.
\begin{itemize}
    \item \textbf{Preparation:} the person lies still while electrodes are placed on the chest, arms, and legs; the skin may be cleaned or shaved for good contact
    \item \textbf{Recording:} the test takes a few minutes and is painless
    \item \textbf{Interpretation:} a doctor or trained clinician reads the trace, often with computer assistance
    \item \textbf{Related tests:} exercise ECGs, 24-hour Holter monitoring for intermittent symptoms, and implantable loop recorders for longer-term monitoring
    \item \textbf{Blood tests:} cardiac enzymes are measured alongside the ECG when a heart attack is suspected
\end{itemize}

\section*{Management}
An ECG is a diagnostic investigation, and management follows from the findings.
\begin{itemize}
    \item \textbf{Normal result:} reassurance and investigation of symptoms through other means if needed
    \item \textbf{Heart attack:} urgent treatment, including medicines and procedures to restore blood flow
    \item \textbf{Arrhythmias:} medicines, cardioversion, pacemakers, or ablation depending on the rhythm
    \item \textbf{Other findings:} management of the underlying condition, with referral to a cardiologist as needed
    \item \textbf{Ongoing monitoring:} repeat ECGs for people with changing symptoms or during treatment
\end{itemize}

\section*{Complications}
\begin{itemize}
    \item The test itself is safe and painless, with no significant complications
    \item Misinterpretation is the main risk, which is why ECGs are read by trained clinicians in context
    \item A false sense of security from a normal resting ECG if symptoms persist
    \item Skin irritation from the electrodes in a small number of people
\end{itemize}

\section*{Prognosis and Outcomes}
The ECG is a fast, reliable tool that saves lives by enabling rapid diagnosis of heart attacks and dangerous arrhythmias. Its value depends on being performed and interpreted correctly and in the right clinical context. When combined with other tests, it allows effective treatment of a wide range of heart conditions.

\section*{Prevention and Self-Care}
\begin{itemize}
    \item Seek prompt medical attention for chest pain, palpitations, or fainting
    \item Attend routine health checks if you have risk factors for heart disease
    \item Carry a copy of previous ECGs or know where they are stored, as comparison is valuable
    \item Follow your doctor's advice on managing blood pressure, cholesterol, and other risk factors
    \item Do not rely on consumer heart monitors for diagnosis; see a doctor for concerns
\end{itemize}

\section*{Sources and Further Reading}
The following reputable sources provide further information for healthcare students and patients seeking reliable, current advice about the ECG.
\begin{itemize}
    \item Heart Foundation Australia. \textit{ECG and heart tests.} https://www.heartfoundation.org.au/
    \item Healthdirect Australia. \textit{Electrocardiogram (ECG).} https://www.healthdirect.gov.au/electrocardiogram
    \item NHS. \textit{Electrocardiogram (ECG): what to expect.} https://www.nhs.uk/conditions/electrocardiogram/
    \item Mayo Clinic. \textit{Electrocardiogram (ECG or EKG).} https://www.mayoclinic.org/tests-procedures/ekg/
    \item MedlinePlus. \textit{Electrocardiogram.} https://medlineplus.gov/ency/article/003868.htm
    \item Better Health Channel. \textit{ECG tests.} https://www.betterhealth.vic.gov.au/
\end{itemize}
